% ------------------------------------------------------------
% Chapter 12: Falsifiability and Discriminating Observations
% ------------------------------------------------------------

\chapter{Falsifiability and Discriminating Observations}
\label{chap:falsifiability}

\section{Overview}

A scientific theory must be falsifiable: it should make predictions that
could, in principle, be shown false by experiment. In this chapter, we
outline potential observational signatures that would falsify \rsvpabbr{},
as well as several predictions that distinguish it from \LCDM{} and
general relativity.

\section{Unique Predictions}

\rsvpabbr{} makes several predictions not shared by \LCDM{}. For example,
the entropic contribution to cosmic acceleration may obviate the need for a
cosmological constant. Specific signatures include deviations in the
distance--redshift relation at high redshift, altered growth rates of
structure, and modifications to CMB anisotropies.

\section{Future Experiments}

Future observations that could test \rsvpabbr{} include next--generation CMB
measurements, galaxy surveys, weak lensing studies, and gravitational wave
observations. Planned missions and observatories offer opportunities to
observe lamphrodynamic effects directly.

\section{Testability Roadmap}

We propose a roadmap identifying key experiments and measurements that could
either support or falsify \rsvpabbr{}. This roadmap emphasizes independent
tests using CMB, LSS, and supernova data, combined with strong--field
observations, to break parameter degeneracies and test core lamphrodynamic
predictions.

\section{Conclusion}

Falsifiability is crucial for scientific progress. While \rsvpabbr{} offers a
promising alternative to \LCDM{}, it must confront observational data
rigorously. Only through careful comparison with current and future
observations can the validity of the lamphrodynamic paradigm be assessed.

\clearpage
